\subsection{Límite uniforme de las medias y suma sobre los niveles}
La regularidad antes de la parada da \(p_i\le C_b/m\), una vacancia diagonal
minorada y una brecha de filas y columnas~\cite[\S7]{pc}.
En consecuencia \(c_b I\le S_0+\Pi_m\le 4I\) para algún \(c_b>0\).
Para ver la cota inferior, la forma cuadrática es la suma de
\(\sum_i p_i(a_i^{\mathsf T}y)^2\) y \(\sum_f h_f y_f^2\), más la proyección de calibre.
El segundo término controla las componentes diagonales y la brecha controla las axiales
ortogonales al calibre. La cota superior sigue de la diagonal unidad de \(S_0\), de que su
masa fuera de la diagonal por fila es a lo sumo tres y de que la proyección actúa en su núcleo.

Sea \(A_m=S_0+\Pi_m-I\), \(M=(I+A_m)^{-1}=R+\Pi_m\) y \(r_b=1-c_b/4<1\).
El polinomio \(P_N=\tfrac14\sum_{h=0}^N[I-(I+A_m)/4]^h\) satisface
\(\|M-P_N\|_{\rm op}\le r_b^{N+1}/c_b\). Para controlar una entrada multiplicada
por \(m\) usamos \(M=I-A_m+A_mMA_m\). La suma absoluta de una fila de \(A_m\) es
a lo sumo cuatro y su máximo es \(O_b(m^{-1})\), luego
\(\|A_me_u\|_2=O_b(m^{-1/2})\) para todo \(u\). Así
\[
 \sup_{u,v}m\,|[A_m(M-P_N)A_m]_{uv}|\le C_b r_b^{N+1}.
\]
Esto acota el error en cada entrada de \(m(R-I)\) al sustituir \(M\) por \(P_N\)
en la identidad anterior. Las entradas originales y sus truncaciones están acotadas.

A \(N\) fijo, las potencias de \(S_0-I\) son sumas sobre sucesiones de líneas adyacentes.
Las inserciones de \(\Pi_m\) se simplifican usando \((S_0-I)\Pi_m=-\Pi_m\).
Un producto de recorridos es una suma finita de tipos de incidencia. Si dos etiquetas afines
son idénticas en un tipo, se cuenta la línea una vez. Una coincidencia adicional impone
una ecuación no constante y deja \(O_{b,N}(n^{d-1})\) posiciones entre las \(O(n^d)\)
de un tipo con \(d\) parámetros. En la diagonal, la ida y vuelta de longitud dos fuerza
la misma casilla en ambos pasos. Su contribución es \(3m\sum_{i\in u}p_i^2\), pues hay tres
elecciones de la otra línea. Su exceso sobre el producto de las medias da precisamente
\(\theta_u=3\eta_u s_u-\int K_{uv}^2\,d\rho(v)\). Los cruces entre diagonales se
cuentan en las clases de paridad de la sección~\ref{sec:contracciones}.

A cada tipo fijo se aplica el seguimiento de líneas hasta el corte auxiliar \(b/4\)
de~\cite[\S4]{tc}. A tolerancia fija \(\eta>0\), fuera de una excepción de probabilidad
exponencial, \(mp_i=W(i)(1+O_b(\eta+m^{-1}))\) en los sitios vivos y
\(h_f=u_f/U_f+O_b(\eta+m^{-1})\). Para un número fijo de marcas el riesgo en un paso es
la suma de sus riesgos individuales salvo \(O_{b,N}(n^{-2})\) por dobles ocupaciones.
Al multiplicar estas supervivencias hasta el nivel \(m\), una marca axial aporta \(x\)
y una diagonal aporta \(U_f(x)\), pues \(\partial_x\log U_f=(1-u_f)/U_f\).
La aproximación es uniforme en las marcas y en \(j\le m\le n\). Se toma primero \(n\to\infty\)
y luego \(\eta\downarrow0\). Las sumas normalizadas de cada clase son sumas de Riemann
sobre politopos fijos. Los factores del perfil son continuos en \(x\) con módulo común en
\([b,1]\) en cada pieza. Por tanto los truncados convergen uniformemente en los niveles
a integrales continuas de los recorridos.

Para aplicar el mismo argumento a \(mV\), sea \(\kappa_i=a_i^{\mathsf T}Ra_i-4\).
La expansión del prefactor~\cite[\S7]{pc} da, con la última suma sobre todas las líneas,
\[
 V=-14\sum_i p_i^3-\tfrac32\sum_i p_i^2\kappa_i
       +\tfrac94\sum_u\left(\sum_{i\in u}p_i^2\right)^2+O_b(m^{-2}).
\]
Las doce entradas fuera de la diagonal en \(m\kappa_i\) aportan \(-12W(i)\) en la
casilla marcada. Combinadas con \(-14\) dan \(4\langle W^2\rangle_i\).
Las partes largas y la diagonal corregida dan los sumandos con \(L_0,\theta\) de
la sección~\ref{sec:profile-data}. El cuadrado de la suma de fibra da su último término.

En las contracciones, el centro tiene masa uno, el exterior masa a lo sumo uno y las
cuatro fibras masa a lo sumo cuatro. La medida de las líneas libres tiene masa \(\le6/b\).
Los pasos de cuadrados tienen masa total \(\le8/b\). Al desarrollar la diferencia de productos,
la cota de cada entrada se multiplica sólo por masas y factores acotados para \(b\) fijo.
Así la cola del grado \(N\) es \(\le C_b r_b^{N+1}\), uniforme en el nivel tanto para la media
física como para su integral límite. El resto de la inserción en la escala \(m^2\) es
\(O_b(m^{-1})\). Tomar primero \(n\to\infty\) y después \(N\to\infty\) da
\begin{align*}
 \sup_{j<m\le n}\Big|
 \mathbb E_\nu\!\left[m^2\sum_{c\in i_m}p_c
       \big(P_c+(\delta_c^0)^2/2\big);H_j\right]-d(m/n)\Big|&\longrightarrow0,\\
 \sup_{j\le m\le n}\left|\mathbb E_\nu[mV(B_m);H_j]-\mathcal V(m/n)\right|
  &\longrightarrow0 .
\end{align*}
Aquí \(i_m\) es la fila revelada y \(H_j\) la llegada regular corta. Para promediar la fila
se usa primero el indicador de supervivencia anterior a su elección. Dado el padre, la fila
es uniforme y la media se convierte en \(m^{-1}\sum_c p_c(\cdot)\).
Cambiar ese indicador por \(\mathbf1_{H_j}\) en los observables acotados cuesta
\(O_b(\nu(H_j^c))\). La segunda fórmula usa todas las familias en el potencial.
Este argumento identifica una media sobre padres y da continuidad del límite en \([b,1]\).
La evaluación decimal de esas integrales por mallas requiere además una cota de su error espacial.

\paragraph{Los extremos y el cambio de ley.}
El error uniforme \(\epsilon_n(b)\) de la primera media contribuye a lo sumo
\(\epsilon_n(b)n\sum_{m>j}m^{-2}=o_b(1)\). El término de \(P_c+(\delta_c^0)^2/2\)
da por tanto \(\int_b^1d(x)x^{-2}\,dx\). Para \(\Delta_cV\), la probabilidad de un hijo
dada la fila bajo \(\nu\) es \(p_c(1+O_b(m^{-1}))\), y \(|\Delta_cV|\le C_bm^{-2}\).
La sustitución en la media cuesta \(O_b(m^{-3})\). Los incrementos centrados hasta la parada
forman una martingala de segundo momento \(O_b(\sum_{m>j}m^{-4})\).
Si \(M_\tau\) es su suma detenida, las diferencias son ortogonales en \(L^2(\nu)\).
Así \(n|\mathbb E_\nu[M_\tau;H_j]|\le n(\mathbb E_\nu M_\tau^2)^{1/2}\nu(H_j)^{1/2}\).
Como \(\sum_{m>j}m^{-4}=O(j^{-3})\), el coste a escala \(n\) es \(O_b(n^{-1/2})\).
En cada historia de \(H_j\) los incrementos realizados suman \(V(B_j)-V(B_n)\).
La segunda media uniforme da \(\mathcal V(b)/b-\mathcal V(1)\).

Hasta la parada, \(n|A|\le C_b\) para el predictor ordinario. Por el Lema~\ref{lem:previa},
\(n|R_Q(n)|\) también está acotado. En \(H_j\) ponemos
\(s=A-R_Q(n)+g_T+\Gamma\). Allí \(|s|=O_b(n^{-1})\) y la densidad inversa de historias
es \(e^{-s-g^\flat}\). Por \eqref{eq:arrival-estimate},
\[
 \int_{H_j}|d\mu-d\nu|
 \le\mathbb E_\mu[|1-e^{-s}|+e^{-s}|1-e^{-g^\flat}|;H_j]=O_b(n^{-1}) .
\]
Con las estimaciones de salida del apéndice~\ref{app:signed} se sigue que
\(n\mathbb E_\mu A=n\mathbb E_\nu[A;H_j]+o_b(1)\).
Para las torres pongamos \(u_e=g_T(M_e^{\rm fis})-g_T(M)\).
La máscara \(M_e^{\rm fis}\) se obtiene del menor simple \(M_e\) aplicando los dos vetos
diagonales del hijo. El cambio es \(O_b(m^{-2})\), incluido el paso que entra en la salida.
Así \(\log\rho^\sharp-\log\rho=\log\mathbb E_{\nu,i}e^{u_e}
 =\mathbb E_{\nu,i}u_e+O_b(m^{-4})\).
La densidad detenida y las cotas ponderadas de salida dan
\(D(\mu^\tau\Vert\nu^\tau)=\mathbb E_\mu[A+\log X]-R_Q(n)=O_b(1)\).
Sean \(d_m\) las entropías condicionales de las dos leyes de hijo dada la fila,
iguales a cero después de la parada. Las filas tienen la misma ley uniforme.
La descomposición de la entropía relativa en sus incrementos condicionales da
\(\sum_m\mathbb E_\mu d_m=D(\mu^\tau\Vert\nu^\tau)\). La desigualdad de
Pinsker~\cite{ct} acota la diferencia de las dos medias de \(u_e\) por \(2C_bm^{-2}\sqrt{d_m/2}\).
Por Jensen y Cauchy--Schwarz, la suma de estas diferencias a escala \(n\) cuesta
\(O_b(n(\sum_{m>j}m^{-4})^{1/2}D(\mu^\tau\Vert\nu^\tau)^{1/2})=O_b(n^{-1/2})\).
Los incrementos realizados de torres se telescopan y añaden
\((F(b)-1/12)/b-(F(1)-1/12)\) a los extremos anteriores. En conclusión, a cada corte fijo,
\[
 \lim_{n\to\infty} n\mathbb E_\mu A^\sharp
 =\int_b^1\frac{d(x)}{x^2}\,dx+\frac{G(b)}b-G(1)
 =\int_b^1\frac{\numerr(x)}{x^2}\,dx .
\]
